\subsection{THE SIMPLE MODULES V(m)}

Let \((u,v)\) be the canonical basis of \(k^2\) . For the identity representation of \(\mathfrak{sl}(2,k)\) ,

\[
X _ {+} u = 0 \quad H u = u \quad X _ {-} u = - v
\]

\[
X _ {+} v = u \quad H v = - v \quad X _ {-} v = 0.
\]

Consider the symmetric algebra \(\mathbf{S}(k^2)\) of \(k^2\) (Algebra, Chap. III, §6, no. 1, Def. 1). The elements of \(\mathfrak{sl}(2,k)\) extend uniquely to derivations of \(\mathbf{S}(k^2)\) , giving \(\mathbf{S}(k^2)\) the structure of an \(\mathfrak{sl}(2,k)\) -module (Chap. I, §3, no. 2). Let \(\mathrm{V}(m)\) be the set of homogeneous elements of \(\mathbf{S}(k^2)\) of degree \(m\) . Then \(\mathrm{V}(m)\) is an \(\mathfrak{sl}(2,k)\) -submodule of \(\mathbf{S}(k^2)\) of dimension \(m + 1\) , the \(m\) th symmetric power of \(\mathrm{V}(1) = k^2\) (Chap. III, Appendix). If \(m,n\) are integers such that \(0 \leq n \leq m\) , put

\[
e _ {n} ^ {(m)} = \binom {m} {n} u ^ {m - n} v ^ {n} \in \mathrm{V} (m).
\]

PROPOSITION 3. For any integer \(m \geq 0\) , \(V(m)\) is an absolutely simple \(\mathfrak{sl}(2, k)\) -module. In this module, \(e_0^{(m)} = u^m\) is primitive of weight \(m\) .

We have \(X_{+}u^{m}=0\) and \(Hu^{m}=mu^{m}\) , so \(u^{m}\) is primitive of weight m. The submodule of \(V(m)\) generated by \(u^{m}\) is of dimension \(m+1\) (Prop. 2(i)) and so is equal to \(V(m)\) . By Prop. 2(v), \(V(m)\) is absolutely simple.

THEOREM 1. Every simple \(\mathfrak{sl}(2,k)\) -module of finite dimension \(n\) is isomorphic to \(\mathrm{V}(n - 1)\) . Every finite dimensional \(\mathfrak{sl}(2,k)\) -module is a direct sum of submodules isomorphic to the modules \(\mathrm{V}(m)\) .

This follows from Lemma 4 and Prop. 1, 2 and 3.

Remarks. 1) The adjoint representation of \(\mathfrak{sl}(2,k)\) defines on \(\mathfrak{sl}(2,k)\) the structure of a simple \(\mathfrak{sl}(2,k)\) -module. This module is isomorphic to V(2) by an isomorphism that takes \(u^2\) to \(X_+\) , \(2uv\) to \(-H\) , and \(v^2\) to \(X_-\) .

\begin{enumerate}
\def\labelenumi{\arabic{enumi})}
\setcounter{enumi}{1}
\tightlist
\item
  For \(n \geq 0\) and m \textgreater{} n,
\end{enumerate}

\[
X _ {-} e _ {n} ^ {(m)} = - (m - n) \binom{m}{n} u ^ {m - n - 1} v ^ {n + 1} = - (n + 1) e _ {n + 1} ^ {(m)}.
\]

Hence, \((e_0^{(m)}, e_1^{(m)}, \ldots, e_m^{(m)})\) is the basis of \(V(m)\) associated to the primitive element \(e_0^{(m)}\) .

\begin{enumerate}
\def\labelenumi{\arabic{enumi})}
\setcounter{enumi}{2}
\tightlist
\item
  Let \(\Phi\) be the bilinear form on \(V(m)\) such that
\end{enumerate}

\[
\begin{array}{c} \varPhi (e _ {n} ^ {(m)}, e _ {n ^ {\prime}} ^ {(m)}) = 0 \text {if} n + n ^ {\prime} \neq m \\ \varPhi (e _ {n} ^ {(m)}, e _ {m - n} ^ {(m)}) = (- 1) ^ {n} \binom {m} {n}. \end{array}
\]

If \(x = au + bv\) and \(y = cu + dv\) , then \(\Phi(x^{m}, y^{m}) = (ad - bc)^{m}\) . It is now easy to check that \(\Phi\) is invariant, and that \(\Phi\) is symmetric for m even, and alternating for m odd.

PROPOSITION 4. Let E be a finite dimensional \(\mathfrak{sl}(2,k)\) -module, \(m\) an integer \(\geq 0\) , \(\mathrm{P}_m\) the set of primitive elements of weight \(m\) . Let L be the vector space of homomorphisms from the \(\mathfrak{sl}(2,k)\) -module \(\mathrm{V}(m)\) to the \(\mathfrak{sl}(2,k)\) -module E. The map \(f \mapsto f(u^m)\) from L to E is linear, injective, and its image is \(\mathrm{P}_m \cup \{0\}\) .

This map is clearly linear, and it is injective because \(u^{m}\) generates the \(\mathfrak{sl}(2,k)\) -module \(\mathrm{V}(m)\) . If \(f \in L\) ,

\[
X _ {+} (f (u ^ {m})) = f (X _ {+} u ^ {m}) = 0, \quad H (f (u ^ {m})) = f (H u ^ {m}) = m f (u ^ {m})
\]

so \(f(u^{m}) \in \mathrm{P}_{m} \cup \{0\}\) . Let \(e \in \mathrm{P}_{m}\) , and \(V\) the submodule of \(E\) generated by \(e\) . By Prop. 1, there exists an isomorphism from the module \(V(m)\) to the module \(V\) that takes \(u^{m}\) to \(e\) . Then \(L(u^{m}) = P_{m} \cup \{0\}\) .

COROLLARY. The isotypical component of E of type V(m) has length

\(\dim (\mathrm{P}_m\cup \{0\} .\)

